\documentclass[11pt,a4paper]{article}
\usepackage{fontspec,amsmath,amssymb,amsthm,microtype}
\setmainfont{lmroman10-regular.otf}[BoldFont=lmroman10-bold.otf,ItalicFont=lmroman10-italic.otf,BoldItalicFont=lmroman10-bolditalic.otf]
\usepackage{xcolor}
\usepackage[most]{tcolorbox}
\usepackage{mathrsfs}
\providecommand{\lt}{<}
\providecommand{\gt}{>}
\setmonofont{DejaVuSansMono.ttf}[BoldFont=DejaVuSansMono-Bold.ttf,ItalicFont=DejaVuSansMono-Oblique.ttf,BoldItalicFont=DejaVuSansMono-BoldOblique.ttf,Scale=0.86]
\newfontfamily\leanfallback{STIXTwoMath-Regular.otf}[Scale=0.86]
\definecolor{leanrule}{gray}{0.62}
\definecolor{leanback}{gray}{0.962}
\definecolor{leantitle}{gray}{0.22}
\definecolor{leanmuted}{gray}{0.38}
\newtcolorbox{leanbox}[1]{enhanced,breakable,colback=leanback,colframe=leanrule,
  boxrule=0.4pt,arc=0.6mm,left=2.4mm,right=2.4mm,top=1.8mm,bottom=1.8mm,
  before skip=0.8em,after skip=1em,fonttitle=\small\bfseries,coltitle=leantitle,
  colbacktitle=leanback,titlerule=0.3pt,title={#1}}
\newlength{\leancharwidth}
\newenvironment{leancode}{\par\smallskip\begingroup\ttfamily\footnotesize\raggedright
  \settowidth{\leancharwidth}{M}\setlength{\parskip}{0pt}\setlength{\parindent}{0pt}
  \frenchspacing}{\par\endgroup\smallskip}
\newcommand{\leanline}[2]{\hangindent=\dimexpr#1\leancharwidth+4\leancharwidth\relax
  \hangafter=1\noindent\hspace*{#1\leancharwidth}#2\par}
\newcommand{\leansource}[1]{\par{\footnotesize\color{leanmuted}#1}\par}
\newcommand{\leanlabel}[1]{\par\smallskip\noindent{\small\bfseries #1}\par}
\usepackage[margin=25mm]{geometry}
\usepackage[unicode,hidelinks]{hyperref}
\hypersetup{pdfauthor={LangevinLean Project}}
\setlength{\emergencystretch}{3em}
\setlength{\parskip}{0.4em}
\setlength{\parindent}{0pt}
\allowdisplaybreaks[2]
\raggedbottom
\title{\texorpdfstring{Real observables and the adjoints of actual evaluations}{Real observables and the adjoints of actual evaluations}}
\author{LangevinLean Project}
\date{}
\begin{document}\maketitle
Draft mathematical exposition; these notes may contain mistakes.


For an actual measure-preserving observation, pointwise real embedding and pointwise adjoint extraction commute with evaluation. Taking adjoints therefore proves that the complex transition compression of an embedded real observable is exactly the embedding of its real transition.


\section{\texorpdfstring{Scalar restriction and pointwise adjoint extraction}{Scalar restriction and pointwise adjoint extraction}}


\textbf{Claim.} Complex adjunction agrees with real adjunction after restricting scalars. On \(L^2\), the real inner product is the integral of the pointwise real inner product and equals the real part of the complex pairing. The adjoint of the lifted embedding is the lift of its pointwise adjoint.



\textbf{Proof.} Test scalar-restricted adjoints against arbitrary vectors and take the real part of the complex adjoint identity. For square-integrable functions, the same identity follows by integrating their pointwise pairings. The real-part functional commutes with the integrable complex pairing; the defining pointwise adjoint equality then identifies the lifted adjoint.

\begin{leanbox}{Lean: \texttt{instance\_\allowbreak{}at\_\allowbreak{}23}, \texttt{instance\_\allowbreak{}at\_\allowbreak{}24}, \texttt{adjoint\_\allowbreak{}restrictScalars} and 4 more}
\leansource{\texttt{instance\_\allowbreak{}at\_\allowbreak{}23} \quad WassersteinMixing/\allowbreak{}ComplexRealTransition.lean, lines 23–23}
\begin{leancode}
\leanline{0}{local\ instance\ :\ InnerProductSpace\ ℝ\ H\ :=\ InnerProductSpace.rclikeToReal\ ℂ\ H}
\end{leancode}
\leansource{\texttt{instance\_\allowbreak{}at\_\allowbreak{}24} \quad WassersteinMixing/\allowbreak{}ComplexRealTransition.lean, lines 24–24}
\begin{leancode}
\leanline{0}{local\ instance\ :\ InnerProductSpace\ ℝ\ K\ :=\ InnerProductSpace.rclikeToReal\ ℂ\ K}
\end{leancode}
\leansource{\texttt{adjoint\_\allowbreak{}restrictScalars} \quad WassersteinMixing/\allowbreak{}ComplexRealTransition.lean, lines 27–34}
\begin{leancode}
\leanline{0}{theorem\ adjoint\_restrictScalars\ (T\ :\ H\ →L[ℂ]\ K)\ :}
\leanline{4}{(T.restrictScalars\ ℝ).adjoint\ =\ T.adjoint.restrictScalars\ ℝ}
\end{leancode}
\leansource{\texttt{instance\_\allowbreak{}at\_\allowbreak{}42} \quad WassersteinMixing/\allowbreak{}ComplexRealTransition.lean, lines 42–42}
\begin{leancode}
\leanline{0}{local\ instance\ :\ InnerProductSpace\ ℝ\ H\ :=\ InnerProductSpace.rclikeToReal\ ℂ\ H}
\end{leancode}
\leansource{\texttt{real\_\allowbreak{}inner\_\allowbreak{}Lp} \quad WassersteinMixing/\allowbreak{}ComplexRealTransition.lean, lines 46–48}
\begin{leancode}
\leanline{0}{theorem\ real\_inner\_Lp\ (μ\ :\ Measure\ A)\ (f\ g\ :\ Lp\ H\ 2\ μ)\ :}
\leanline{4}{inner\ ℝ\ f\ g\ =\ ∫\ a,\ inner\ ℝ\ (f\ a)\ (g\ a)\ ∂μ}
\end{leancode}
\leansource{\texttt{real\_\allowbreak{}inner\_\allowbreak{}Lp\_\allowbreak{}eq\_\allowbreak{}re} \quad WassersteinMixing/\allowbreak{}ComplexRealTransition.lean, lines 52–55}
\begin{leancode}
\leanline{0}{theorem\ real\_inner\_Lp\_eq\_re\ (μ\ :\ Measure\ A)\ (f\ g\ :\ Lp\ H\ 2\ μ)\ :}
\leanline{4}{inner\ ℝ\ f\ g\ =\ (inner\ ℂ\ f\ g).re}
\end{leancode}
\leansource{\texttt{mapLp\_\allowbreak{}adjoint} \quad WassersteinMixing/\allowbreak{}ComplexRealTransition.lean, lines 58–77}
\begin{leancode}
\leanline{0}{theorem\ mapLp\_adjoint\ (i\ :\ R\ →ₗᵢ[ℝ]\ H)\ (μ\ :\ Measure\ A)\ :}
\leanline{4}{(mapLp\ i\ μ).toContinuousLinearMap.adjoint\ =}
\leanline{6}{i.toContinuousLinearMap.adjoint.compLpL\ 2\ μ}
\end{leancode}
\leanlabel{How the Lean proof follows the argument}
The local scalar restrictions specify the underlying real inner products. \texttt{adjoint\_\allowbreak{}restrict\allowbreak{}Scalars} uses uniqueness of adjoints. \texttt{real\_\allowbreak{}inner\_\allowbreak{}Lp\_\allowbreak{}eq\_\allowbreak{}re} uses the integral of the real part. \texttt{map\allowbreak{}Lp\_\allowbreak{}adjoint} tests against arbitrary real observables and compares exact almost-everywhere representatives.
\end{leanbox}

\section{\texorpdfstring{Naturality of actual evaluation and its adjoint}{Naturality of actual evaluation and its adjoint}}


\textbf{Claim.} If \(\phi_*\nu=\mu\), then pointwise embedding and extraction commute with pullback by \(\phi\). The adjoint pullback satisfies the corresponding embedding identity.



\textbf{Proof.} Both sides have the same almost-everywhere representative, obtained by composing the pointwise map with \(\phi\). A measure-preserving map transports the exceptional null sets correctly. Express extraction as the proved lifted adjoint and take adjoints of its commuting diagram. Scalar-restriction compatibility and the double-adjoint identity give the asserted adjoint-evaluation diagram.

\begin{leanbox}{Lean: \texttt{extraction\_\allowbreak{}compMeasurePreserving}, \texttt{mapLp\_\allowbreak{}compMeasurePreserving}, \texttt{evaluation\_\allowbreak{}adjoint\_\allowbreak{}mapLp}}
\leansource{\texttt{extraction\_\allowbreak{}compMeasurePreserving} \quad WassersteinMixing/\allowbreak{}ComplexRealTransition.lean, lines 80–94}
\begin{leancode}
\leanline{0}{theorem\ extraction\_compMeasurePreserving\ (i\ :\ R\ →ₗᵢ[ℝ]\ H)}
\leanline{4}{(μ\ :\ Measure\ A)\ (ν\ :\ Measure\ B)\ (φ\ :\ B\ →\ A)\ (hφ\ :\ MeasurePreserving\ φ\ ν\ μ)}
\leanline{4}{(f\ :\ Lp\ H\ 2\ μ)\ :}
\leanline{4}{i.toContinuousLinearMap.adjoint.compLpL\ 2\ ν\ (Lp.compMeasurePreserving\ φ\ hφ\ f)\ =}
\leanline{6}{Lp.compMeasurePreserving\ φ\ hφ\ (i.toContinuousLinearMap.adjoint.compLpL\ 2\ μ\ f)}
\end{leancode}
\leansource{\texttt{mapLp\_\allowbreak{}compMeasurePreserving} \quad WassersteinMixing/\allowbreak{}ComplexRealTransition.lean, lines 98–109}
\begin{leancode}
\leanline{0}{theorem\ mapLp\_compMeasurePreserving\ (i\ :\ R\ →ₗᵢ[ℝ]\ H)}
\leanline{4}{(μ\ :\ Measure\ A)\ (ν\ :\ Measure\ B)\ (φ\ :\ B\ →\ A)\ (hφ\ :\ MeasurePreserving\ φ\ ν\ μ)}
\leanline{4}{(f\ :\ Lp\ R\ 2\ μ)\ :}
\leanline{4}{mapLp\ i\ ν\ (Lp.compMeasurePreserving\ φ\ hφ\ f)\ =}
\leanline{6}{Lp.compMeasurePreserving\ φ\ hφ\ (mapLp\ i\ μ\ f)}
\end{leancode}
\leansource{\texttt{evaluation\_\allowbreak{}adjoint\_\allowbreak{}mapLp} \quad WassersteinMixing/\allowbreak{}ComplexRealTransition.lean, lines 112–131}
\begin{leancode}
\leanline{0}{theorem\ evaluation\_adjoint\_mapLp\ (i\ :\ R\ →ₗᵢ[ℝ]\ H)}
\leanline{4}{(μ\ :\ Measure\ A)\ (ν\ :\ Measure\ B)\ (φ\ :\ B\ →\ A)\ (hφ\ :\ MeasurePreserving\ φ\ ν\ μ)}
\leanline{4}{(f\ :\ Lp\ R\ 2\ ν)\ :}
\leanline{4}{(Lp.compMeasurePreservingₗᵢ\ ℂ\ φ\ hφ).toContinuousLinearMap.adjoint\ (mapLp\ i\ ν\ f)\ =}
\leanline{6}{mapLp\ i\ μ\ ((Lp.compMeasurePreservingₗᵢ\ ℝ\ φ\ hφ).toContinuousLinearMap.adjoint\ f)}
\end{leancode}
\leanlabel{How the Lean proof follows the argument}
The two composition theorems use \texttt{Lp.\allowbreak{}ext} and transport representative equalities through the measure-preserving map. \texttt{evaluation\_\allowbreak{}adjoint\_\allowbreak{}map\allowbreak{}Lp} takes adjoints of the exact extraction diagram and applies \texttt{adjoint\_\allowbreak{}restrict\allowbreak{}Scalars}.
\end{leanbox}

\section{\texorpdfstring{Actual transition compression preserves the real image}{Actual transition compression preserves the real image}}


\textbf{Claim.} Let \(U_0,U_1\) be evaluation isometries of two observations with the same marginal. Their complex compression \(U_0^*U_1\) sends the embedded real observable \(if\) to \(i(U_{0,\mathbb R}^*U_{1,\mathbb R}f)\), hence preserves the actual real image.



\textbf{Proof.} Move embedding through the second evaluation using pointwise naturality, then through the first adjoint using the previous theorem. The resulting vector is explicitly the embedding of the genuine real compression, which supplies its witness in the real range.

\begin{leanbox}{Lean: \texttt{compression\_\allowbreak{}mapLp}, \texttt{compression\_\allowbreak{}mem\_\allowbreak{}range}}
\leansource{\texttt{compression\_\allowbreak{}mapLp} \quad WassersteinMixing/\allowbreak{}ComplexRealTransition.lean, lines 135–146}
\begin{leancode}
\leanline{0}{theorem\ compression\_mapLp\ (i\ :\ R\ →ₗᵢ[ℝ]\ H)}
\leanline{4}{(μ\ :\ Measure\ A)\ (ν\ :\ Measure\ B)\ (φ\ ψ\ :\ B\ →\ A)}
\leanline{4}{(hφ\ :\ MeasurePreserving\ φ\ ν\ μ)\ (hψ\ :\ MeasurePreserving\ ψ\ ν\ μ)}
\leanline{4}{(f\ :\ Lp\ R\ 2\ μ)\ :}
\leanline{4}{ObservableTransition.compression\ (Lp.compMeasurePreservingₗᵢ\ ℂ\ φ\ hφ)}
\leanline{6}{(Lp.compMeasurePreservingₗᵢ\ ℂ\ ψ\ hψ)\ (mapLp\ i\ μ\ f)\ =}
\leanline{6}{mapLp\ i\ μ\ (ObservableTransition.compression\ (Lp.compMeasurePreservingₗᵢ\ ℝ\ φ\ hφ)}
\leanline{8}{(Lp.compMeasurePreservingₗᵢ\ ℝ\ ψ\ hψ)\ f)}
\end{leancode}
\leansource{\texttt{compression\_\allowbreak{}mem\_\allowbreak{}range} \quad WassersteinMixing/\allowbreak{}ComplexRealTransition.lean, lines 149–155}
\begin{leancode}
\leanline{0}{theorem\ compression\_mem\_range\ (i\ :\ R\ →ₗᵢ[ℝ]\ H)}
\leanline{4}{(μ\ :\ Measure\ A)\ (ν\ :\ Measure\ B)\ (φ\ ψ\ :\ B\ →\ A)}
\leanline{4}{(hφ\ :\ MeasurePreserving\ φ\ ν\ μ)\ (hψ\ :\ MeasurePreserving\ ψ\ ν\ μ)}
\leanline{4}{(f\ :\ Lp\ R\ 2\ μ)\ :}
\leanline{4}{ObservableTransition.compression\ (Lp.compMeasurePreservingₗᵢ\ ℂ\ φ\ hφ)}
\leanline{6}{(Lp.compMeasurePreservingₗᵢ\ ℂ\ ψ\ hψ)\ (mapLp\ i\ μ\ f)\ ∈\ Set.range\ (mapLp\ i\ μ)}
\end{leancode}
\leanlabel{How the Lean proof follows the argument}
\texttt{compression\_\allowbreak{}map\allowbreak{}Lp} composes the evaluation and adjoint-evaluation identities. \texttt{compression\_\allowbreak{}mem\_\allowbreak{}range} exhibits the actual real compression as the preimage, so range preservation is proved from the observation maps rather than introduced as an assumption.
\end{leanbox}

\end{document}
